\section{\label{II_2_C}Processus de validation technique}

%Poursuite de l'application du concept HIL 

Afin de gérer les essais-erreurs de l'affinage du modèle sur les données spécifiques du corpus, il est pertinent de mettre en place un processus de validation technique du traitement. Ce processus permettrait d'une part de valider les résultats du modèle, et d’autre part, d'assurer la concordance entre les données produites par le modèle avec les données déjà existantes au sein du logiciel.

La validation des prédictions du modèle peut être effectuée au sein de la plateforme \gls{labels}. Les archives brittaniques proposent que le jeu de données test utilisé précédemment afin de corriger les prédictions, soit distinct d'un second jeu de données test, dédié à l'évaluation des performances du modèle, à différents niveaux de supervision humaine \footcite{kasturi2024}. Sur ce second jeu de données, sont notamment intégrées des métadonnées de supervision humaine. En effet, Kasturi explique l'introduction des notations "HIL-10 \%, HIL-15 \% et HIL-20 \%, qui signifient respectivement que les utilisateurs ont corrigé 10 \%, 15 \% et 20 \% des boîtes englobantes prédites peu performantes dans l’ensemble de données de test." \footcite{kasturi2024}. (traduction de :  HIL-10\%, HIL-15\%, and HIL-20\%, which signify that users have corrected 10\%, 15\%, and 20\% of poorly performing predicted bounding boxes in the test dataset, respectively). Les données corrigées manuellement remplacent ensuite certaines images au hasard dans le premier jeu de données d'apprentissage, contituant un nouveau jeu de données d'entraînement sur lequel le modèle s'affine \footcite{kasturi2024}. 
Le modèle ainsi affiné est finalement exécuté sur le second jeu de données, corrigé selon différents niveaux d'intervention humaine \footcite{kasturi2024}. Pour chaque niveau de supervision, sont analysés les "scores moyens d’Intersection over Union (mIOU) et de précision moyenne (mAP), ainsi que de la précision et du rappel" \footcite{kasturi2024}. 
Les résultats obtenus montrent que les performances du modèle s'améliorent considérablement à mesure que le pourcentage d'images corrigées manuellement augmente \footcite{kasturi2024}. En effet, l'enrichissement du jeu de données permet au modèle d'appréhender l'analyse des images de manière plus précise. L'auteur note toutefois la nécessité de poursuivre l'évaluation du modèle sur d'autres données afin de vérifier son adaptabilité et ses performances pour la détection d'autres types d'objets \footcite{kasturi2024}. En effet, le modèle devrait être ré entraîné constamment, à mesure qu'il soit exécuté sur de nouvelles données. (Active learning)


De plus, le processus de validation technique implique également la validation de l'indexation des données produites au sein du logiciel, en adéquation avec les règles métier en place. Au sein de l’espace de travail S-Movies dédié au corpus étudié, l'indexation des mots-clés est contrôlée par le thésaurus "keyword". Or, le modèle \gls{yolo} est nativement pré entraîné sur un jeu de données généraliste dont les étiquettes produites ne correspondront pas directement aux autorités indexées dans la base de données. Deux cas peuvent être gérés par l'utilisateur côté technique. D'une part, les mots-clés produits par le modèle peuvent déjà être indexés sous une forme contrôlée au sein du logiciel. Dans ce cas, il revient à l'utilisateur de vérifier que les mots-clés générés par le modèle concordent avec la forme indexée contrôlée au sein du logiciel. À terme, cette tâche de validation pourra être gérée par un LLM, un grand modèle de langage, déjà employé par Skinsoft pour un outil de recherche en langage naturel par exemple. Ce grand modèle de langage pourra être capable, à partir du thésaurus indexé dans le logiciel, de valider les formes contrôlées des mots-clés générés par \gls{yolo}. Pour cela, un mapping (tableau de corespondance) pourra être fourni au LLM entre les mots-clés que le modèle peut potentiellement générer, et leur forme contrôlée dans le thésaurus. Nous avons effectivement déjà élaboré un tel tableau au cours de notre stage dans le cadre de la mise en place de l'intelligence artificielle pour d'autres tâches, notamment la mise en place d'une recherche en langage naturel. Pour cette tâche tout à fait différente que celle que nous décrivons dans ce mémoire, nous avons effectué une table de correspondance entre les champs contrôlés pouvant être requêtés dans le logiciel et leur définition en langage naturel, préparant ainsi la gestion de synonymes et de termes employés par l'utilisateur potentiellement pour requêter ces champs. Sur le même principe, nous pourrions mettre en place une table de correspondance afin d'assurer la validation des données produites et d'assurer \textit{in fine} leur indexation au sein de l'application. Un calcul de similarité aide le LLM à repérer les termes simiaires et els rapprocher de leur forme contrôlée.
D'autre part, si les mots clés n'existent pas sous une forme contrôlée au sein du thésaurus et qu'aucune correspondance n'est détectée, un nouveau terme de thésaurus pourra être créé au sein du logiciel, avec l'accord de l'utilisateur.

%à ajouter : Named Entity Linking – NEL compare les entités trouvées avec les données normatives). L'algorithme ou le modèle mathématique effectue une classification multi-labels. Il se base sur des concepts statistiques tels que TF/IDF (Term Frequency/Inverted Document Frequency). \footcite{zotero-item-584}